\documentclass[10pt,a4paper]{amsart}
\usepackage[margin=19mm]{geometry}
\usepackage{amsmath,amssymb,mathtools}
\usepackage[scr=rsfs,cal=euler]{mathalfa}
\usepackage{xfrac}
\usepackage[unicode,hidelinks,bookmarks=false]{hyperref}
\newcommand{\ie}{i.e.\ }
\newcommand{\cf}{cf.\ }
\newcommand{\resp}{respectively\ }
\newcommand{\Subsection}{\subsection}
\providecommand{\nobreakdash}{\nobreak\hbox{-}}
\setlength{\emergencystretch}{2em}
\multlinegap0pt
\usepackage{bm}
\usepackage{bbm}
\usepackage{dsfont}
\usepackage{xspace}
\usepackage{cancel}
\usepackage{mathtools}
\usepackage{amsthm}
\makeatletter
\@ifundefined{theorem}{\newtheorem{theorem}{Theorem}}{}
\@ifundefined{definition}{\newtheorem{definition}[theorem]{Definition}}{}
\@ifundefined{lemma}{\newtheorem{lemma}[theorem]{Lemma}}{}
\@ifundefined{proposition}{\newtheorem{proposition}[theorem]{Proposition}}{}
\makeatother
\def\A{\mathcal{A}}
\def\Crit{\operatorname{Crit}}
\def\End{\operatorname{End}}
\def\Gr{\operatorname{Gr}}
\def\Hess{\operatorname{Hess}}
\def\Hom{\operatorname{Hom}}
\def\Id{\operatorname{Id}}
\def\R{\mathbb{R}}
\def\Tr{\operatorname{Tr}}
\def\U{\operatorname{U}}
\def\ad{\operatorname{ad}}
\def\grad{\operatorname{grad}}
\def\pr{\operatorname{pr}}
\def\so{\mathfrak{so}}
\def\sp{\mathfrak{sp}}
\def\u{\mathfrak{u}}
\title[A Morse-Bott unification of the Grassmannians of a symplecti]{A Morse-Bott unification of the Grassmannians of a symplectic vector space}
\author{Hyunmoon Kim}
\date{Source: arXiv:2601.16441v2 (CC BY 4.0)}
\begin{document}
\maketitle

\noindent\emph{Source and licence.} A Morse-Bott unification of the Grassmannians of a symplectic vector space. Version arXiv:2601.16441v2 (CC BY 4.0), distributed under the Creative Commons Attribution 4.0 International licence, as recorded in the arXiv licence metadata for this exact version and shown on the abstract page https://arxiv.org/abs/2601.16441v2. Full rights notice: licenses/arxiv-2601.16441-CC-BY-4.0.txt.\par\smallskip

\noindent\textit{Editorial scope.} Counted unit: the Theorem whose source label is \texttt{thm:morse\_bott} in the article (source0.tex lines 810--812, source label \texttt{thm:morse\_bott}), delivered together with the article's own complete original proof of it (source0.tex lines 814--903, one \texttt{proof} environment of 90 lines). No mathematics is written, repaired or re-derived: statement and proof are the article's own text line for line. Required context retained uncounted: lem:commutators (lines 578--594); def:energy (lines 632--641); lem:derivatives (lines 700--711); prop:critical\_locus (lines 754--757). Every definition, display and label of those lines is retained unchanged. Omitted from this package and not redistributed: 1--577, 595--631, 642--699, 712--753, 758--809, 813--813, 904--1118. Nothing inside a retained excerpt is omitted or altered: every retained body is delivered exactly as the article's lines are written, so no omission receipt of any kind accompanies this package. Citations: the retained excerpts carry no citation command at all, so no bibliography entry of the article is transcribed and no reference is redistributed. Reference adaptation: none was needed and none was made; every reference inside the delivered unit resolves inside it, so no unresolved reference or citation is produced. Printed slips kept literal: no printed word, symbol, hypothesis or number of the counted statement or of its proof was altered. Layout and class: the article's own class is \texttt{article} with packages including fullpage, amsthm, amsmath, amssymb, amsfonts, tikz-cd, mathtools, comment, authblk, biblatex, draftwatermark, showlabels; the wrapper is the lane's pinned \texttt{amsart} editorial wrapper at 10pt on A4, which restates the theorem-like environments the retained text needs and adds the article's own macro definitions that the retained text needs (\texttt{\textbackslash A}, \texttt{\textbackslash Crit}, \texttt{\textbackslash End}, \texttt{\textbackslash Gr}, \texttt{\textbackslash Hess}, \texttt{\textbackslash Hom}, \texttt{\textbackslash Id}, \texttt{\textbackslash R}, \texttt{\textbackslash Tr}, \texttt{\textbackslash U}, \texttt{\textbackslash ad}, \texttt{\textbackslash grad}, \texttt{\textbackslash pr}, \texttt{\textbackslash so}, \texttt{\textbackslash sp}, \texttt{\textbackslash u}), with extra wrapper packages (bm, bbm, dsfont, xspace, cancel, mathtools). The delivered numbering is the unit's own. Assets: no retained span contains a figure environment, a graphics-inclusion command, a PDF-page inclusion, a table or a TikZ picture; no image file is copied into this package, the package declares no asset dependency and no image file is redistributed with it. Licence: version arXiv:2601.16441v2 of this article is distributed under the Creative Commons Attribution 4.0 International licence, as recorded in the arXiv licence metadata for this exact version and shown on the abstract page https://arxiv.org/abs/2601.16441v2. A full rights notice is delivered at licenses/arxiv-2601.16441-CC-BY-4.0.txt. Characters: every retained excerpt is pure ASCII, so the delivered engine, XeLaTeX with its default Latin Modern text font, reports no missing character.
\begin{lemma}\label{lem:commutators}\
    \begin{enumerate}
        \item There exists a short exact sequence
        \[ 0 \to \so(W) \times \so(W^\perp) \to \so(V) \xrightarrow{-\ad_{P_W}} T_{W}\Gr(k; V) \to 0.\]
        Moreover, for any $\xi \in \so(V)$ its fundamental vector field $X_\xi \in \Gamma (T\Gr(k; V))$ is given by
        \[ (X_\xi)_P:= -[P, \xi].\]
        \item Let $\psi_\cdot : \so(V) \to \End (\Gamma(T\Sigma|_{\Gr(k; V)}))$ be defined by
        \[ \psi_\xi(X)_P:= [\widetilde{\xi}_P, X_P]\]
        where $\widetilde{\xi}$ is the constant section of the trivial vector bundle $\Gr(k; V) \times \so(V)$ with value $\xi$. Then $\psi$ is a Lie algebra representation of $\so(V)$.
         \item The orthogonal projection $\pr_P: T_P\Sigma\twoheadrightarrow T_P\Gr(k; V)$ is given by
        \[ \pr_P(A) = \big[P, [P, A]\big].\]
        \item The shape operator $\A_\cdot: T_P\Sigma|_{\Gr(k; V)} \to \End(T_P\Gr(k; V))$ is given by
        \[ \A_X (Y)_P = \big[P, [Y, X]\big] \quad X \in T_P\Sigma|_{\Gr(k; V)}, Y \in T_P\Gr(k; V).\]
        If $X \in T_P \Gr(k; V) \subseteq T_P\Sigma|_{\Gr(k; V)}$ then $\A_X(Y)_P = 0$.
       
    \end{enumerate}
\end{lemma}

\begin{definition} \label{def:energy}
Let $X_J \in \Gamma(T\Gr(k; V))$ be the vector field defined by
\[ (X_J)_P:= [P, J] = \ad_{P}(J).\]

Let the \textbf{energy function} $f: \Gr(k; V) \to \R$ be defined as follows:
\[
f(P) := \frac{1}{2} \lVert (X_J)_P \rVert^2 = \frac{1}{2} \Tr\left( [P, J] [P, J]^T \right) = \frac{1}{2} \Tr([P, J]^2)
\]
where for $W \in \Gr(k; V)$, $P = P_W$ is the orthogonal projection  operator to $W$.
\end{definition}

\begin{lemma}\label{lem:derivatives} \
Let $f: \Gr(k; V) \to \R$ be the energy function defined in Definition \ref{def:energy}.
\begin{enumerate}
    \item The differential of $f$ evaluated at $Y$ is the trace pairing with a double commutator:
    \[ (df)_P(Y) = \langle \psi_J(X_J)_P, Y\rangle =   \Tr\left( \big[ J, [P, J]\big] Y \right). \]
    \item The gradient vector field of $f$ is given by:
    \[ (\grad f)_P = \pr_P (\psi_J(X_J)_P) = \big[ P, [P, [J, [P, J]]] \big]. \]
    \item At any critical point $P \in \Crit(f)$, the Hessian operator $\Hess(f)_P:T_P\Gr(k; V) 
    \to T_P\Gr(k; V)$ is given by:
    \begin{align*}\Hess(f)_P(Y) &= -\pr_P (\psi_J^2(Y)) + \A_{\psi_J(X_J)_P}(Y) \\&= \bigg[P, \big[P, \big [J, [Y, J]\big]\big]\bigg] + \bigg[P, \big[Y, \big[J, [P, J]\big] \big] \bigg].\end{align*}
\end{enumerate}
\end{lemma}

\begin{proposition}[Critical Locus] \label{prop:critical_locus}
The critical locus of the energy function $f: \Gr(k; V) \to \R$ is  the union the Grassmannians of $J$-compatible linear subspaces:
\[ \Crit(f) = \bigsqcup_{\vec{n}: n_0 + 2n_+ = k} \Gr_J(\vec{n}; V).\]
\end{proposition}

\begin{theorem}[Morse-Bott Property]\label{thm:morse_bott}
The energy function $f: \Gr(k; V) \to \R$ is a Morse-Bott function. Its critical set is the disjoint union of the smooth submanifolds $\Gr_J(\vec{n}; V)$, and its Hessian is non-degenerate on the normal bundle of these submanifolds.
\end{theorem}

\begin{proof}
The first claim is proved in Proposition~\ref{prop:critical_locus}. 

Fix $P=P_W\in \Gr_J(\vec n;V)$. Since $W$ is $J$-compatible, we have orthogonal decompositions
\[
W=W_0\oplus W_+^J,\qquad W^\perp=JW_0\oplus W_-^J.
\]

\smallskip
Choose an orthonormal basis adapted to the ordered splitting
\[
V = W_0\oplus W_+^J\oplus JW_0\oplus W_-^J.
\]
In this basis,
\[
P=
\begin{pmatrix}
\Id_{n_0} & 0 & 0 & 0\\
0 & \Id_{2n_+} & 0 & 0\\
0 & 0 & 0 & 0\\
0 & 0 & 0 & 0
\end{pmatrix},
\qquad
J=
\begin{pmatrix}
0 & 0 & -\Id_{n_0} & 0\\
0 & J_+ & 0 & 0\\
\Id_{n_0} & 0 & 0 & 0\\
0 & 0 & 0 & J_-
\end{pmatrix},
\]
where $J_\pm^2=-\Id$ and $J_\pm^T=-J_\pm$. 
\smallskip

Using the identification
\begin{align*}
    T_P\Gr(k; V) &= \Hom(W, W^\perp) = \Hom(W_0 \oplus W_+^J, W_-^J\oplus JW_0)\\ &\cong \Hom(W_0, JW_0) \oplus \Hom(W_+^J, JW_0) \oplus \Hom(W_0, \oplus W_-^J) \oplus \Hom(W_+^J, W_-^J),
\end{align*}
we can express a tangent vector $Y \in T_P\Gr(k; V)$ as
\[ Y = (A, B, C, D)\]
where $A\in\Hom(W_0,JW_0)$, $B\in\Hom(W_+^J,JW_0)$, $C\in\Hom(W_0,W_-^J)$, $D\in\Hom(W_+^J,W_-^J)$.
Relative to the splitting $V=W_0\oplus W_+^J\oplus JW_0\oplus W_-^J$,
\[
Y=
\begin{pmatrix}
0 & 0 & A^T & C^T\\
0 & 0 & B^T & D^T\\ 
A & B & 0 & 0\\
C & D & 0 & 0
\end{pmatrix}
\]

\smallskip

In this notation we can compute as in Lemma \ref{lem:derivatives}
\begin{align*}
 -\pr_P(\psi_J^2(Y)) &= [P, [P, [J, [Y, J]]]] = (2(A + A^T), 2B, 2C, 2D + 2J_- D J_+)\\
\A_{\psi_J(X_J)_P}(Y) &= [P, [Y, [J, [P, J]]]] = (-4A, -2B, -2C, 0).\end{align*}

Summing yields the Hessian action:
\begin{equation}\label{eq:hessian-action-correct}
\Hess(f)_P(Y) =  \big(2(A^T-A),\,0,\,0,\,2(D+J_-DJ_+)\big).
\end{equation}

Hence
\begin{equation}\label{eq:ker-hess-correct}
\ker \Hess(f)_P=\{(A,B,C,D):\ A=A^T,\ \ J_-D=DJ_+\},
\end{equation}
and the normal directions are exactly the skew part of $A$ and the $J$-antilinear part of $D$; in particular the Hessian is nondegenerate on the normal bundle.

\smallskip
The critical submanifold containing $P$ is the $\U(V,J)$-orbit passing through $P$, so by Lemma~\ref{lem:commutators}(1) applied to the $\U(V,J)$-action
\[
T_P\Gr_J(\vec n;V)=\ad_P(\u(V,J))=\{-[P,\xi]:\ \xi\in\u(V,J)\}.
\]

The elements of the Lie algebra $\u(V, J) = \sp(V) \cap \so(V)$ are matrices $\xi$ such that $\xi^TJ + J\xi = 0$ and $\xi = -\xi^T$. Equivalently, they are skew-symmetric matrices $X$ that commute with $J$. With respect to the splitting $V = W_0\oplus W_+^J\oplus JW_0\oplus W_-^J$, they are
\[ \xi = \begin{pmatrix} E & BJ_+ & -A^T & -C^T \\ -J_+ B & F & - B^T & -D^T \\ A & B & E & -C^TJ_- \\ C & D & J_-C & G \end{pmatrix} \]
such that the entries satisfy the following:
\begin{itemize}
    \item $E$, $F$, $G$ are skew-symmetric, and $[F, J_+] = 0$ and $[G, J_-] = 0$, i.e.
\[ (E, F, G) \in \so(W_0) \times \u(W_+^J, J_+) \times \u(W_-^J, J_-)\]
\item $A$ is symmetric, and $D$ is complex linear $J_-D = DJ_+$.
\end{itemize}
Computing the commutator, we obtain
\[ [\xi, P] = \begin{pmatrix} 0 & 0 & -A^T & -C^T \\ 0 & 0 & - B^T & -D^T \\ A & B & 0 & 0 \\ C & D & 0 & 0 \end{pmatrix}.\]
and thus we have
\[ \ad_P(\u(V, J)) = T_P\Gr_J(\vec{n}; V) = \ker \Hess(f)_P,\]
which is the Morse-Bott condition.
\end{proof}

\end{document}
